\documentclass{article}
\usepackage{amsmath}
\usepackage{amssymb}
\usepackage{graphicx}
\usepackage{hyperref}
\usepackage{parskip}

\title{A Persistence Framework for Life, Intelligence, and Science}
\author{Albert Jan van Hoek}
\date{\today}

\begin{document}

\maketitle

\section*{Introduction}
The Fermi paradox is usually asked as: \textit{where is everybody?} But the deeper question may be: \textit{what actually persists?}

\section*{Life as Persistence}
Life is not a thing. It is a process that holds itself together against a world that does not care whether it continues. Its defining feature is not complexity or intelligence. It is continuation under conditions that are always changing. Life is what refuses to stop.

But life did not invent this rule. It inherited it.

Everything we observe --- every particle, every structure, every system --- is a survivor. The universe is not a neutral backdrop against which persistence occasionally appears. It is the accumulated residue of everything that passed the filter. What we see when we look out is not emptiness punctuated by rare complexity. It is proof, everywhere, that the same principle holds: what remains is what found a way to remain.

\section*{Reframing the Fermi Paradox}
This reframes the Fermi paradox entirely. The silence is not a mystery. It is the answer. The universe looks exactly as it should if only persistent things leave traces.

\section*{Intelligence and Science as Persistence Strategies}
Intelligence fits inside this picture --- but it does not crown it.

Complexity achieves efficiency gains in persistence through many routes. A crystal finds stability through geometry. An ecosystem finds it through redundancy and feedback. An immune system finds it through memory and adaptation. Intelligence is one more solution to the same problem: how to continue under conditions that change faster than simple structure can track.

Intelligence is not special in kind. It is remarkable in degree. It is what happens when a persistent process begins modeling its own conditions --- when persistence becomes aware of what persistence requires.

Science is what happens when that awareness becomes collective and self-correcting.

We tend to tell a particular story about science: that it is humanity's restless desire to know, to explore, to push into the unknown. That story is not wrong. But it is incomplete --- and the part it leaves out may be the most important part.

Science is also --- perhaps primarily --- how persistence scales. Not one organism modeling its environment, but many minds building a shared, correctable picture of what must remain stable for the whole system to continue. It is the mechanism by which a process that began with the first self-replicating molecule eventually learns what it needs to survive the civilization it built.

\section*{The Purpose of Science}
This changes what science is for.

Discovery is the method. Persistence is the function. We go out to understand the unknown because the unknown contains the conditions we depend on --- and we cannot protect what we cannot see. Every time science makes a hidden dependency visible --- a tipping point in a climate, a threshold in an immune response, a feedback loop in an economy --- it is not just expanding knowledge. It is buying time. It is the process, looking ahead, keeping itself going.

A civilization that treats science only as curiosity --- as a luxury of prosperous times --- has misunderstood what it is holding. It is not a monument to human ingenuity. It is the infrastructure of continuation.

\section*{Conclusion}
Which means we are not observers of this system. We are not even its most advanced expression.

We are the moment one branch of it turns, looks at the whole, and recognizes the rule it has been running on since the beginning.

What persists is whatever learns what persistence requires --- in time.

The Fermi silence is not a question about where intelligence went. It is a reminder that intelligence was never the point. Persistence was always the point. Intelligence is just one of its more talkative solutions.

\end{document}